% =========================================================
% CHAPTER 3 — PROBLEM FORMULATION AND OBJECTIVES
% =========================================================

\chapter{Problem Formulation and Objectives}

\section{Problem Statement}

The contemporary digital landscape is heavily characterized by an increasing asymmetry in information and control between service providers and end-users. While web technologies have advanced to offer dynamic, highly personalized user experiences, they have simultaneously enabled sophisticated mechanisms for user manipulation and exploitation. This manipulation predominantly manifests in two distinct but related forms: dark patterns and phishing attacks. 

Dark patterns represent a deliberate subversion of human-computer interaction (HCI) principles. Rather than optimizing for usability and transparency, these user interfaces are meticulously engineered to exploit human cognitive biases, thereby nudging or coercing users into actions that are disproportionately beneficial to the service provider. Examples include forced continuity in subscription models, hidden costs that only appear at the final stage of a checkout process, and confirmshaming that guilt-trips users into opting into data collection. The fundamental problem is that conventional security mechanisms—such as web firewalls and antivirus software—treat dark patterns as benign content because they are structurally identical to legitimate UI components and are hosted on legitimate domains.

Conversely, phishing attacks represent a more overt security threat wherein malicious entities masquerade as trustworthy entities to extract sensitive credentials, such as banking passwords and personal identification information. Traditional anti-phishing mechanisms rely heavily on static blacklists, URL reputation scoring, and domain age analysis. However, modern phishing campaigns utilize zero-day domains, perfectly cloned user interfaces, and contextually persuasive language, rendering static heuristics ineffective.

A critical vulnerability in modern AI-based threat detection systems is the high incidence of false positives on legitimate authentication forms. Legitimate sign-in and credential verification interfaces (such as GitHub, Google, LinkedIn, and PayPal) heavily feature critical security keywords like "login", "password", "verify", "secure", and "account". In standard sequence classification models, the presence of these keywords triggers high-confidence phishing classifications (>0.95), even when the domain is fully trusted and legitimate. This lack of domain-aware context leads to alert fatigue, prompting users to disable security controls and exposing them to real attacks.

The core problem addressed by this project is the lack of a unified, real-time, context-aware detection system capable of identifying both manipulative interface designs (dark patterns) and fraudulent semantic intents (phishing) directly within the user's browsing environment without falsely flagging trusted resources. Existing solutions are either heavily siloed—focusing solely on phishing while ignoring dark patterns—or rely on post-hoc analysis rather than real-time intervention. There is a critical need for an intelligent client-side agent that can continuously parse the Document Object Model (DOM), verify domain reputations, analyze text heuristics, and run transformer-based inference asynchronously to highlight malicious elements while keeping legitimate sites clean.

\section{Objectives of the Project}

The primary objective of this project is to conceptualize, design, and implement an AI-powered browser extension, named ``DESIGN TO DECEIVE'', capable of detecting and mitigating dark patterns and phishing attempts in real-time. To achieve this overarching goal, the project is divided into the following specific objectives:

\begin{enumerate}
    \item \textbf{To develop a dynamic DOM scanning architecture:} Design a lightweight content script capable of traversing the web page's Document Object Model in real-time, utilizing \texttt{MutationObserver} APIs to handle dynamically loaded content in Single Page Applications (SPAs) without causing significant main-thread blocking or rendering latency.
    \item \textbf{To implement a robust NLP-based classification backend:} Engineer a high-performance RESTful API using FastAPI to serve fine-tuned RoBERTa (Robustly Optimized BERT Pretraining Approach) transformer models. This backend must be capable of receiving batched text payloads from the extension and performing binary or multi-class semantic classification with minimal inference delay.
    \item \textbf{To construct distinct detection pipelines for multifaceted threats:} Train and deploy specialized NLP models for different threat vectors: one optimized for identifying coercive and manipulative language indicative of dark patterns, and another fine-tuned on phishing corpora to detect deceptive social engineering tactics.
    \item \textbf{To engineer an intuitive, non-disruptive user interface for threat visualization:} Develop a visual highlighting mechanism and a contextual tooltip system that overlays onto the web page, alerting users to specific deceptive elements. Additionally, create a centralized cybersecurity dashboard within the extension popup to display aggregated risk scores, confidence metrics, and historical detection data.
    \item \textbf{To implement a hybrid, context-aware decision engine:} Design a decision pipeline that integrates URL reputation checks, domain whitelisting, urgency keyword matching, and transformer inference to differentiate between trusted authentication pages and malicious credential harvesting attempts.
    \item \textbf{To ensure real-time performance and privacy preservation:} Optimize the communication overhead between the browser extension and the backend inference server using asynchronous message passing and payload batching, ensuring that the detection pipeline operates within acceptable latency thresholds suitable for active browsing.
\end{enumerate}

% =========================================================
% CHAPTER 4 — SYSTEM ARCHITECTURE AND METHODOLOGY
% =========================================================

\chapter{System Architecture and Methodology}

\section{Proposed Methodology}

The proposed methodology for the ``DESIGN TO DECEIVE'' framework relies on a hybrid architecture combining client-side DOM extraction with server-side heavy-duty NLP inference. The methodology is structured around a continuous, real-time processing loop that executes asynchronously alongside the user's browsing activities.

Unlike traditional static analysis tools that evaluate an entire webpage's source code at load time, our approach adopts a dynamic, element-wise analysis strategy. As the user navigates a website, the extension selectively targets specific HTML elements that are historically prone to containing manipulative text—such as buttons, hyperlinks, form labels, and modal dialogs. The text content and the current page URL are extracted, cleaned of HTML markup, and transmitted to the centralized machine learning backend.

At the backend, the system implements a hybrid decision logic. First, the page domain is extracted and verified against a trusted whitelist. If the domain is trusted, the classification thresholds are tightened and standard input fields are bypassed to prevent false alerts. Second, the domain is checked against URL reputation rules (such as suspicious TLDs, excessive hyphens, and token-based brand spoofing). Third, the text is scanned for contextual triggers like urgency language. Finally, the text is processed by fine-tuned RoBERTa sequence classification models to establish a contextual threat signature.

The classification result, along with its confidence score, risk tier (None, Low Risk, Medium Risk, High Risk), and a generated explanation, is returned to the extension. The extension then modifies the local DOM to inject visual indicators—specifically, CSS-based highlights and interactive tooltips—exclusively for High Risk detections, preventing interface clutter on benign pages.

\section{System Architecture}

The system architecture follows a decoupled, client-server paradigm, specifically chosen to offload the computationally expensive transformer inference from the resource-constrained browser environment to a dedicated backend server.

The architecture comprises three primary tiers:

\begin{enumerate}
    \item \textbf{The Client Tier (Chrome Extension):} Operating directly within the user's browser, this tier consists of Content Scripts (for DOM manipulation and extraction), a Background Service Worker (acting as the central event router and state manager), and the User Interface components (Popup Dashboard and injected Tooltips).
    \item \textbf{The Middleware Tier (FastAPI Engine):} A high-throughput Python-based web server responsible for handling incoming HTTP requests, managing asynchronous batching, and routing payloads to the appropriate NLP models and heuristic modules.
    \item \textbf{The Inference Tier (Hybrid Engine and RoBERTa Models):} The core intelligence of the system, combining pre-trained language models loaded into memory via PyTorch, a domain reputation evaluator, and a heuristic rule compiler to calculate contextual threat bands.
\end{enumerate}

\section{Working Principle}

The working principle of the system can be delineated into five sequential phases:

\textbf{Phase 1: Initialization and Injection.} Upon navigating to a new URL, the browser executes the extension's content script. The script registers event listeners and initializes a \texttt{MutationObserver} to monitor the DOM for structural changes.

\textbf{Phase 2: Extraction and Filtering.} The content script scans the DOM for targeted tags (e.g., \texttt{<a>}, \texttt{<button>}, \texttt{<span>}). It extracts the text, the tag name, and the page URL, generating a local batch of elements with unique identifiers.

\textbf{Phase 3: Payload Transmission.} The content script utilizes the Chrome Messaging API to pass the batched elements to the background service worker, which forwards an asynchronous HTTP \texttt{POST} request to the FastAPI backend.

\textbf{Phase 4: Hybrid Inference and Logging.} The backend receives the payload and processes each element. The domain is parsed and reputation heuristics are run. The backend checks if the domain is whitelisted. If the domain is trusted and the text contains no urgency indicators, the backend skips the ML model pass, logging the bypass. Otherwise, the text is tokenized and evaluated by the RoBERTa model. The final threat level is computed by integrating the ML logits with domain and urgency heuristics. Detections are assigned to risk bands (None, Low, Medium, High). High-confidence alerts are logged as security warnings.

\textbf{Phase 5: Visualization and Alerting.} The service worker receives the response. The content script processes each item: elements categorized as High Risk phishing are visually outlined with a pulsating red CSS border, and mouse hover listeners are registered to render dynamic tooltips. Medium and Low risk classifications are cached for reporting but do not visually alter the webpage layout.

\section{Data Flow Diagram}

To illustrate the flow of data through the system, Data Flow Diagrams (DFDs) have been designed at various abstraction levels.

The Level 0 DFD provides a macroscopic view, showing the user initiating web requests, the browser fetching HTML content, the extension intercepting the DOM along with page URLs, and the interaction with the external NLP API.

The Level 1 DFD breaks down the system into detailed processes: DOM Extraction, Message Routing, HTTP Request Handling, Domain Whitelist Check, URL Reputation Scoring, RoBERTa Tokenization and Inference, Risk Score Normalization, and DOM Modification.
